%==============================================================
% ABSTRAK DAN ABSTRACT
%==============================================================

\clearpage
\thispagestyle{empty}

% --- ABSTRAK (Bahasa Indonesia) ---
\begin{center}
{\fontsize{14pt}{17pt}\selectfont\bfseries ABSTRAK}
\end{center}

\vspace{6pt}

\noindent
GHIFFARI BRAVIA HISHAM, MUHAMMAD ABYAN PUTRA WIBOWO, ADAM NAUFAL,
MOCHAMAD CHAIRULRIDJAL NURVIKRI, RAFIF MUHAMMAD FARRAS.
Klasifikasi \textit{Stage Retinopathy of Prematurity} Berdasarkan Citra Fundus Bayi
Menggunakan \textit{Enhancement} Citra, Segmentasi Pembuluh Darah, dan \textit{Custom} CNN.
Mata kuliah Pengenalan Citra Digital. Dibimbing oleh NAMA DOSEN.

\vspace{12pt}

\noindent
\textit{Retinopathy of Prematurity} (ROP) adalah gangguan perkembangan pembuluh darah retina
pada bayi prematur yang berpotensi menyebabkan kebutaan jika tidak terdeteksi sejak dini.
Penelitian ini membangun sebuah \textit{pipeline} klasifikasi \textit{stage} ROP dari citra
fundus bayi ke lima kelas: Normal, Stage 1, Stage 2, Stage 3, dan bekas luka laser
(\textit{laser scars}). Proses dilakukan dalam dua tahap utama, yaitu segmentasi pembuluh darah
secara klasik dan klasifikasi berbasis jaringan saraf tiruan. Pipeline segmentasi
menggabungkan filter Gabor multiskala dengan respons Meijering berbasis matriks Hessian
(konfigurasi \textit{g40\_m60\_fine}), menghasilkan nilai Dice terbaik 0,4739 pada dataset
HVDROPDB-BV. Model klasifikasi utama adalah \textit{TinyResNetV2} yang dilatih dari awal tanpa
bobot pra-latih, menerima masukan tiga kanal berupa \textit{softmap} pembuluh darah, \textit{softmap} area tepi, dan kanal hijau ternormalisasi CLAHE. Validasi silang 5-\textit{fold} yang
mempertimbangkan kelompok citra serupa (\textit{group-aware cross-validation}) menghasilkan
\textit{macro F1-score} sebesar \textbf{0,7853}. Percobaan ablasi yang terkontrol menunjukkan bahwa
representasi \textit{softmap} memberikan peningkatan 0,0987 dibanding masukan RGB mentah
(0,6866 vs 0,7853) pada model, data, dan protokol yang sama. Hasil ini membuktikan bahwa
representasi spasial dari informasi pembuluh darah---bukan hanya nilai piksel mentah---merupakan
kunci utama peningkatan performa klasifikasi \textit{stage} ROP.

\vspace{12pt}

\noindent
Kata kunci: \textit{custom} CNN, \textit{enhancement} citra, citra fundus,
\textit{Retinopathy of Prematurity}, segmentasi pembuluh darah, \textit{TinyResNetV2}

\vspace{1.5cm}

% --- ABSTRACT (English) ---
\begin{center}
{\fontsize{14pt}{17pt}\selectfont\bfseries ABSTRACT}
\end{center}

\vspace{6pt}

\noindent
GHIFFARI BRAVIA HISHAM, MUHAMMAD ABYAN PUTRA WIBOWO, ADAM NAUFAL,
MOCHAMAD CHAIRULRIDJAL NURVIKRI, RAFIF MUHAMMAD FARRAS.
Classification of Retinopathy of Prematurity Stages from Infant Fundus Images
Using Image Enhancement, Vessel Segmentation, and Custom CNN.
Digital Image Recognition Course. Supervised by NAMA DOSEN.

\vspace{12pt}

\noindent
Retinopathy of Prematurity (ROP) is a vascular developmental disorder of the retina in
premature infants that can lead to blindness if not detected early. This study constructs a
classification pipeline for ROP staging from infant fundus images into five classes: Normal,
Stage 1, Stage 2, Stage 3, and laser scars. The process consists of two main stages: classical
vessel segmentation and neural network-based classification. The segmentation pipeline combines
multiscale Gabor filters with Hessian-based Meijering response (the \textit{g40\_m60\_fine} configuration),
achieving a best Dice score of 0.4739 on the HVDROPDB-BV dataset. The main classification model
is TinyResNetV2 trained from scratch without pretrained weights, receiving a three-channel input
consisting of vessel softmap, ridge softmap, and CLAHE-normalized green channel. Group-aware
5-fold cross-validation yields a macro F1-score of \textbf{0.7853}. A controlled ablation experiment
demonstrates that the softmap representation provides a +0.0987 improvement over raw RGB input
(0.6866 vs. 0.7853) under identical model, data, and protocol settings. These results confirm
that spatial vessel evidence---rather than raw pixel values---is the key driver of classification
performance improvement for ROP staging.

\vspace{12pt}

\noindent
Keywords: custom CNN, image enhancement, fundus image, Retinopathy of Prematurity,
vessel segmentation, TinyResNetV2
